% Setting Up, Chapter 4 -- Your Worksheet Files (annex Sec 1: the SU part's
% chapters on the disks and files the program uses).
% Sources (Law 1): os/i123/app/face_a.c file_name (the name rules: no
% space, * or ?; the default .wk1 extension; a drive or folder may be
% typed), paint_command (the prompts and the default A:\*.wk1, the current
% file offered by /File Save), retrieve / save / write_sheet (the
% transaction; the 64K file buffer; Cancel / Replace / Backup), os/i123/fs/
% wk1.c (the .WK1 codec and THE CARRY: records the program does not use
% are written back unchanged; test-i123-wk1-roundtrip over the 44 goldens).
% GATES: test-i123-wk1-roundtrip test-flair-i123-files test-flair-i123-file-ui test-pubs-i123
\Gates{test-i123-wk1-roundtrip, test-flair-i123-files,
  test-flair-i123-file-ui, test-pubs-i123}
\chapter{Your Worksheet Files}

A worksheet you build exists only in the computer's memory until you save
it in a file.  This chapter tells you where Initech 123 keeps your
worksheet files, how to name them, and how to look after them.

\section{Where Your Files Are Kept}
\ix{data disk}
Initech 123 saves and retrieves worksheet files on the InitechOS data disk,
drive \file{A:}.  When it asks for a file name, the control panel shows
\texttt{A:\textbackslash *.wk1}: a file name you type without a drive or a
folder is a file in the main folder of the data disk, beside the folder
\file{APPS} that holds Initech 123 itself.

\section{Naming Your Files}
\ix{file names}
A file name has up to eight letters and digits, optionally followed by a period
and an extension of up to three characters.  You can type it in capital or
small letters.  A file name cannot contain a space or the characters
\texttt{*} and \texttt{?}.

Initech 123 gives every worksheet file the extension \file{.WK1} unless you
type another one, and looks for that extension when you retrieve a file.
Choose names that say what the worksheet holds:
\file{BUDGET92}, \file{SALESQ3}, \file{LOAN}.

\section{The Current File}
\ix{current file}
After you save or retrieve a worksheet, its file is the \term{current
file}.  The next time you select \slashcmd{/File Save}, Initech 123 offers
the name of the current file, so that you need only press \kname{ENTER}
and then select Replace to save the worksheet again.  The status line
shows \texttt{Saved} or \texttt{Retrieved} and the name of the file after
each save or retrieve.

\section{Looking After Your Files}
\ix{backup}
\begin{itemize}
\item Save your work often, and always before you erase the worksheet or
  retrieve another one: each of these replaces the worksheet in memory.
\item To keep an earlier version of a worksheet, save the changed worksheet
  under a new name, such as \file{BUDGET2}.  Saving under the old name
  and selecting Replace overwrites the earlier version.
\item Keep copies of your worksheet files on diskettes, as your InitechOS
  books describe for any file of the data disk.
\end{itemize}

\section{Worksheet Files from Other Programs}
\ix{.WK1 files}
Initech 123 reads and writes worksheet files in the \file{.WK1} format.
A \file{.WK1} worksheet file made by another program can be retrieved if
it is not larger than 64K and does not hold more entries than Initech 123
can (Appendix~A of the \doctitle{Initech 123 Reference}).  Settings in the
file that Initech 123 does not use, such as print and graph settings, are
kept and written back unchanged when you save the worksheet.

If a file cannot be retrieved, Initech 123 tells you why on the control
panel, and the worksheet you were working on stays as it was.
